\chapter{External Distribution Sort, Over-Sampling and Random Sampling}
\label{chap:external-distribution-sort-en}

\begin{center}
  \textbf{Lecture 6}\\
  \emph{External Distribution Sort, Over-Sampling with Chernoff Bounds, and Random Sampling}\\[0.5em]
  Date: September 30, 2026
\end{center}

In the context of hierarchical memory models and out-of-core sorting (\textit{External Memory Model}), the \textbf{Multiway Merge Sort} studied in \autoref{cha:lecture-three} represents the canonical \textit{bottom-up} paradigm: data elements are initially arranged into sorted sequences (\textit{runs}) and subsequently merged recursively in $k$-way passes until a single sorted array is obtained.

However, a dual paradigm of fundamental theoretical and practical significance exists: \textbf{External Distribution Sort} (often referred to as out-of-core \textit{Multiway Quicksort}). This approach embraces a \textit{top-down} philosophy: instead of merging pre-sorted runs, it partitions the entire unsorted dataset into multiple disjoint partitions (\textit{buckets}), each corresponding to an ordered interval of key values, and recursively processes each bucket independently.

In this chapter, we investigate the architecture of External Distribution Sort, addressing the critical challenge of balanced pivot selection through the statistical technique of \textbf{Over-Sampling}. We then formally derive the failure probability bound using \textbf{Chernoff Bounds} (measure concentration inequalities) and examine uniform \textbf{Random Sampling} algorithms, comparing hash-based rejection methods with the streaming paradigm of \textbf{Reservoir Sampling}.

\section{External Distribution Sort and \texorpdfstring{$k$}{k}-Way Partitioning}
\label{sec:external-distribution-sort-model-en}

Consider a dataset comprising $n$ elements stored on secondary storage under the two-level memory model of Aggarwal and Vitter~\cite{aggarwal1988}, characterized by an internal primary memory capacity of $M$ elements and block transfers of size $B$ elements. The goal is to sort the array $S[1 \dots n]$ while minimizing the total number of I/O operations.

\subsection{Operational Mechanism and Bucket Definition}
The algorithm partitions the array $S$ into $k$ contiguous, disjoint partitions, termed \textbf{buckets} $B_1, B_2, \dots, B_k$, through the selection of $k - 1$ strictly increasing \textbf{pivots}:
\begin{equation}
  s_1 < s_2 < \dots < s_{k-1}.
\end{equation}

To handle boundary conditions uniformly, two conceptual sentinel values are introduced:
\begin{equation}
  s_0 \coloneqq -\infty, \qquad s_k \coloneqq +\infty.
\end{equation}

\begin{definition}[Bucket $B_i$]
For each index $i \in \{1, \dots, k\}$, bucket $B_i$ is the subset of elements from $S$ whose keys fall within the half-open interval bounded by adjacent pivots:
\begin{equation}
  B_i \coloneqq \left\{ x \in S \mid s_{i-1} < x \le s_i \right\}.
  \label{eq:bucket-definition-en}
\end{equation}
\end{definition}

The topology of $k$-way partitioning is illustrated in \autoref{fig:k-way-partitioning-model-en}.

\begin{figure}[H]
  \centering
  \resizebox{0.95\textwidth}{!}{%
  \begin{tikzpicture}[
    scale=1.0,
    bucket/.style={rectangle, draw=blue!80!black, thick, minimum height=1.3cm, align=center, font=\small\bfseries},
    sentinel/.style={font=\footnotesize\bfseries, text=red!80!black},
    pivot/.style={font=\footnotesize\bfseries, text=blue!85!black},
    ptr/.style={-{Stealth[scale=0.9]}, very thick, draw=blue!85!black}
  ]
    % Bucket segments
    \node[bucket, fill=blue!12, minimum width=2.8cm, anchor=west] (b1) at (0, 0) {
      $B_1$\\[2pt]\footnotesize ($s_0 < x \le s_1$)
    };
    \node[bucket, fill=orange!15, minimum width=2.8cm, anchor=west] (b2) at (b1.east) {
      $B_2$\\[2pt]\footnotesize ($s_1 < x \le s_2$)
    };
    \node[bucket, fill=teal!15, minimum width=2.8cm, anchor=west] (b3) at (b2.east) {
      $B_3$\\[2pt]\footnotesize ($s_2 < x \le s_3$)
    };
    \node[draw=none, minimum width=1.5cm, anchor=west, align=center, font=\Large\bfseries] (dots) at (b3.east) {
      $\dots$
    };
    \node[bucket, fill=purple!15, minimum width=2.8cm, anchor=west] (bk) at (dots.east) {
      $B_k$\\[2pt]\footnotesize ($s_{k-1} < x \le s_k$)
    };

    % Pivots and Sentinels at top
    \node[above=14pt, sentinel] at (b1.north west) {$s_0 = -\infty$};
    \draw[ptr] ([yshift=18pt]b1.north east) -- (b1.north east)
      node[pos=0.0, above, pivot] {$s_1$};
    \draw[ptr] ([yshift=18pt]b2.north east) -- (b2.north east)
      node[pos=0.0, above, pivot] {$s_2$};
    \draw[ptr] ([yshift=18pt]b3.north east) -- (b3.north east)
      node[pos=0.0, above, pivot] {$s_3$};
    \draw[ptr] ([yshift=18pt]bk.north west) -- (bk.north west)
      node[pos=0.0, above, pivot] {$s_{k-1}$};
    \node[above=14pt, sentinel] at (bk.north east) {$s_k = +\infty$};

    % Decorative bottom braces
    \draw[decorate, decoration={brace, amplitude=5pt, mirror}, thick, blue!85!black]
      ([yshift=-4pt]b1.south west) -- ([yshift=-4pt]b1.south east)
      node[midway, below=7pt, font=\scriptsize\bfseries] {$|B_1| \approx \frac{n}{k}$};

    \draw[decorate, decoration={brace, amplitude=5pt, mirror}, thick, orange!90!black]
      ([yshift=-4pt]b2.south west) -- ([yshift=-4pt]b2.south east)
      node[midway, below=7pt, font=\scriptsize\bfseries] {$|B_2| \approx \frac{n}{k}$};

    \draw[decorate, decoration={brace, amplitude=5pt, mirror}, thick, teal!90!black]
      ([yshift=-4pt]b3.south west) -- ([yshift=-4pt]b3.south east)
      node[midway, below=7pt, font=\scriptsize\bfseries] {$|B_3| \approx \frac{n}{k}$};

    \draw[decorate, decoration={brace, amplitude=5pt, mirror}, thick, purple!85!black]
      ([yshift=-4pt]bk.south west) -- ([yshift=-4pt]bk.south east)
      node[midway, below=7pt, font=\scriptsize\bfseries] {$|B_k| \approx \frac{n}{k}$};
  \end{tikzpicture}%
  }
  \caption{Key universe partitioning into $k$ disjoint buckets using $k-1$ sorted pivots and two conceptual sentinel boundaries.}
  \label{fig:k-way-partitioning-model-en}
\end{figure}

\begin{property}[Immediate Concatenation Property]
Unlike Merge Sort, where the final phase necessitates merging sorted streams, External Distribution Sort satisfies:
\begin{equation}
  \forall x \in B_i, \quad \forall y \in B_j \quad \text{with } i < j \implies x \le y.
\end{equation}
Consequently, once each individual bucket $B_i$ has been recursively sorted, the globally sorted array is simply the direct sequential concatenation of all buckets:
\begin{equation}
  S_{\text{sorted}} = B_1 \circ B_2 \circ \dots \circ B_k.
\end{equation}
No subsequent merging step is required, eliminating pointer overhead and post-processing computations entirely.
\end{property}

\subsection{The Balancing Objective and Its Impact on I/O}
For the algorithm to achieve optimal I/O complexity in the external memory model, the choice of branching factor $k$ and pivot locations must satisfy two architectural constraints:

\begin{enumerate}
  \item \textbf{Internal buffer capacity:} During the distribution pass over $S$, internal RAM must maintain at least one block-sized write buffer of size $B$ for each of the $k$ buckets. Since total memory is $M$, the maximum admissible partitioning degree is:
  \begin{equation}
    k = \Theta\left(\frac{M}{B}\right), \quad \text{with } k \ge 2.
    \label{eq:k-branching-factor-en}
  \end{equation}
  Whenever the write buffer for bucket $B_i$ becomes full, it is flushed to disk via a single sequential block write of $B$ elements.
  
  \item \textbf{Dimensional bucket balance:} Each bucket must be evenly sized:
  \begin{equation}
    |B_i| = \Theta\left(\frac{n}{k}\right), \quad \forall i \in \{1, \dots, k\}.
    \label{eq:balance-condition-en}
  \end{equation}
\end{enumerate}

\begin{theorem}[I/O Complexity of Balanced External Distribution Sort]
When balance condition~\eqref{eq:balance-condition-en} is maintained across all recursive levels, the depth of the recursion tree satisfies:
\begin{equation}
  h = \mathcal{O}\left(\log_k \frac{n}{M}\right) = \mathcal{O}\left(\log_{M/B} \frac{n}{B}\right).
\end{equation}
Because each recursion level performs a linear distribution pass over the entire dataset costing $\mathcal{O}(n/B)$ block transfers, the overall I/O complexity is:
\begin{equation}
  \text{I/O}(n) = \mathcal{O}\left(\frac{n}{B} \log_{M/B} \frac{n}{B}\right),
\end{equation}
matching the theoretical sorting lower bound established by Aggarwal and Vitter.
\end{theorem}

\begin{property}[Catastrophic Degeneracy from Skewed Buckets]
If even a single bucket $B_j$ is oversized (e.g., $|B_j| = \Omega(n)$), tree height ceases to shrink logarithmically and degenerates linearly toward $\Omega(n/M)$. Under this pathological case, total I/O would collapse to quadratic complexity $\mathcal{O}((n/B)^2)$. Accurate pivot selection is therefore the central pillar of the entire algorithm.
\end{property}

\section{The Over-Sampling Technique}
\label{sec:oversampling-technique-en}

How can we identify $k - 1$ pivots that guarantee, with overwhelming confidence, that each of the $k$ buckets contains $\Theta(n/k)$ elements?

A naive approach would simply sample $k - 1$ random elements from $S$ uniformly and sort them. However, order statistics theory for uniform random variables shows that the maximum bucket size exhibits high variance, with an expected maximum bucket size of:
\begin{equation}
  \mathbb{E}\left[\max_{1 \le i \le k} |B_i|\right] \approx \frac{n}{k} \ln k.
\end{equation}
For realistic values of $k$ ($k = M/B \approx 10^3 \div 10^5$), the factor $\ln k \approx 7 \div 12$ introduces unacceptable skew.

To resolve this limitation without computing exact medians and quantiles on disk (which would incur heavy linear I/O overheads), the \textbf{Over-Sampling} technique is deployed.

\subsection{Parameters and Operational Procedure}
We introduce an oversampling multiplier $a = \Theta(\log_2 k)$, termed the \textbf{oversampling factor}. As derived analytically below, we set:
\begin{equation}
  a \coloneqq 12 \ln k.
  \label{eq:oversampling-factor-en}
\end{equation}

The pivot selection procedure operates in three concise stages:

\begin{enumerate}
  \item \textbf{Global sample extraction:} Draw $(a + 1)k - 1$ items uniformly at random with replacement from $S$, accumulating them into an auxiliary internal buffer $A$:
  \begin{equation}
    |A| = (a + 1)k - 1.
    \label{eq:sample-size-en}
  \end{equation}
  
  \item \textbf{Internal sorting of sample:} Sort array $A$ entirely within fast internal memory using an optimal comparison sort. Since $|A| = \mathcal{O}(k \ln k) = \mathcal{O}(\frac{M}{B} \ln \frac{M}{B}) \ll M$, sorting $A$ requires 0 disk I/O and negligible CPU time.
  
  \item \textbf{Equidistant pivot selection:} Select the $k - 1$ pivots by picking elements from sorted array $A$ at regular strides of $(a + 1)$:
  \begin{equation}
    s_i \coloneqq A[(a + 1) \cdot i], \quad \text{for } i = 1, \dots, k - 1.
    \label{eq:pivot-selection-rule-en}
  \end{equation}
\end{enumerate}

The internal structure of sorted sample $A$ is illustrated in \autoref{fig:oversampling-sample-layout-en}.

\begin{figure}[H]
  \centering
  \resizebox{0.95\textwidth}{!}{%
  \begin{tikzpicture}[
    scale=1.0,
    box/.style={rectangle, draw=black!80, thick, minimum height=1.1cm, align=center, font=\small},
    pivotBox/.style={rectangle, draw=blue!90!black, fill=blue!20, thick, minimum height=1.1cm, minimum width=1.1cm, align=center, font=\small\bfseries},
    arr/.style={-{Stealth[scale=0.9]}, very thick, draw=blue!85!black}
  ]
    % Cells of sample A
    \node[box, fill=gray!12, minimum width=2.4cm, anchor=west] (gap0) at (0, 0) {$a$ samples};
    \node[pivotBox, anchor=west] (p1) at (gap0.east) {$s_1$};
    
    \node[box, fill=gray!12, minimum width=2.4cm, anchor=west] (gap1) at (p1.east) {$a$ samples};
    \node[pivotBox, anchor=west] (p2) at (gap1.east) {$s_2$};
    
    \node[box, fill=gray!12, minimum width=2.4cm, anchor=west] (gap2) at (p2.east) {$a$ samples};
    \node[draw=none, minimum width=1.2cm, anchor=west, font=\large\bfseries] (dots) at (gap2.east) {$\dots$};
    \node[pivotBox, anchor=west] (pk) at (dots.east) {$s_{k-1}$};
    
    \node[box, fill=gray!12, minimum width=2.4cm, anchor=west] (gapk) at (pk.east) {$a$ samples};

    % Index pointers above
    \draw[arr] ([yshift=16pt]p1.north) -- (p1.north)
      node[pos=0.0, above, font=\scriptsize\bfseries, text=blue!90!black] {$A[a+1]$};
      
    \draw[arr] ([yshift=16pt]p2.north) -- (p2.north)
      node[pos=0.0, above, font=\scriptsize\bfseries, text=blue!90!black] {$A[2(a+1)]$};

    \draw[arr] ([yshift=16pt]pk.north) -- (pk.north)
      node[pos=0.0, above, font=\scriptsize\bfseries, text=blue!90!black] {$A[(k-1)(a+1)]$};

    % Overall bottom brace
    \draw[decorate, decoration={brace, amplitude=6pt, mirror}, thick]
      ([yshift=-4pt]gap0.south west) -- ([yshift=-4pt]gapk.south east)
      node[midway, below=8pt, font=\small\bfseries] {Overall sorted sample $|A| = (a+1)k - 1$};
  \end{tikzpicture}%
  }
  \caption{Layout of sorted auxiliary array $A$: exactly $a$ intermediate sample items lie between any two consecutive pivots.}
  \label{fig:oversampling-sample-layout-en}
\end{figure}

\begin{property}[Intermediate Sample Invariant]
By algorithmic design, between any two consecutive pivots $s_{i-1}$ and $s_i$, there reside \textbf{exactly $a$ intermediate samples} from $A$. Across all $k$ intervals, this accounts for $k \cdot a$ elements; adding the $k - 1$ pivots yields:
\begin{equation}
  k \cdot a + (k - 1) = (a + 1)k - 1 = |A|,
\end{equation}
confirming perfect conservation of sample capacity.
\end{property}

\subsection{Concrete Numerical Trace}
\begin{example}[Oversampling Trace with $n=64$, $k=4$, $a=3$]
To make the mechanics tangible, let $n = 64$ and target $k = 4$ buckets with $a = 3$:
\begin{enumerate}
  \item Sample size:
  \begin{equation}
    |A| = (a + 1)k - 1 = (3 + 1) \cdot 4 - 1 = 15 \text{ elements}.
  \end{equation}
  \item Suppose uniform random sampling and internal sorting produce:
  \begin{equation*}
    A = [\, 4,\, 7,\, 9,\, \mathbf{12},\, 15,\, 18,\, 22,\, \mathbf{29},\, 33,\, 38,\, 41,\, \mathbf{50},\, 55,\, 61,\, 68 \,].
  \end{equation*}
  \item Selection stride is $a + 1 = 4$. The $k - 1 = 3$ pivots are:
  \begin{align*}
    s_1 &= A[1 \cdot 4] = A[4] = \mathbf{12}, \\
    s_2 &= A[2 \cdot 4] = A[8] = \mathbf{29}, \\
    s_3 &= A[3 \cdot 4] = A[12] = \mathbf{50}.
  \end{align*}
  \item The 4 resulting bucket boundaries are:
  \begin{align*}
    B_1 &= \{ x \in S \mid -\infty < x \le 12 \}, \\
    B_2 &= \{ x \in S \mid 12 < x \le 29 \}, \\
    B_3 &= \{ x \in S \mid 29 < x \le 50 \}, \\
    B_4 &= \{ x \in S \mid 50 < x \le +\infty \}.
  \end{align*}
\end{enumerate}
Each bucket interval $(s_{i-1}, s_i]$ encloses exactly $a = 3$ items from sample $A$, faithfully tracking the empirical key density.
\end{example}

\section{Failure Probability Bound via Chernoff Bounds}
\label{sec:chernoff-bound-proof-en}

We now formally prove that oversampling guarantees dimensional bucket balance with overwhelming probability.

\subsection{Formal Definition of Failure}
Pivot extraction results in a \textbf{failure} if at least one bucket $B_i$ is severely imbalanced, holding at least four times its expected average capacity $n/k$:
\begin{equation}
  \mathbb{P}(\text{failure}) \coloneqq \mathbb{P}\left(\exists i \in \{1, \dots, k\} : |B_i| \ge \frac{4n}{k}\right).
  \label{eq:failure-definition-en}
\end{equation}

\subsection{Geometric Construction on Conceptual Sorted Array \texorpdfstring{$S_{\text{sort}}$}{S\_sort}}
Consider the hypothetically sorted array $S_{\text{sort}}$. Subdivide $S_{\text{sort}}$ into $\frac{k}{2}$ contiguous segments of uniform width, denoted $t_j$ for $j = 1, \dots, \frac{k}{2}$:
\begin{equation}
  |t_j| \coloneqq \frac{n}{k/2} = \frac{2n}{k}.
  \label{eq:segment-width-en}
\end{equation}

The geometric relationship between segment $t_j$ and an oversized bucket $B_i$ is visualized in \autoref{fig:pigeonhole-chernoff-construction-en}.

\begin{figure}[H]
  \centering
  \resizebox{0.95\textwidth}{!}{%
  \begin{tikzpicture}[
    scale=1.0,
    segment/.style={rectangle, draw=black!80, thick, minimum height=1.1cm, align=center, font=\small},
    bucketBox/.style={rectangle, draw=red!80!black, fill=red!15, very thick, minimum height=1.2cm, align=center, font=\small\bfseries}
  ]
    % Sorted array S_sort partitioned into k/2 segments of size 2n/k
    \node[segment, fill=gray!10, minimum width=2.4cm, anchor=west] (t1) at (0, 0) {$t_1$\\\footnotesize($\frac{2n}{k}$)};
    \node[segment, fill=yellow!20, minimum width=2.4cm, anchor=west] (t2) at (t1.east) {$t_2$\\\footnotesize($\frac{2n}{k}$)};
    \node[segment, fill=orange!25, minimum width=2.4cm, anchor=west] (t3) at (t2.east) {$t^* = t_3$\\\footnotesize($\frac{2n}{k}$)};
    \node[segment, fill=yellow!20, minimum width=2.4cm, anchor=west] (t4) at (t3.east) {$t_4$\\\footnotesize($\frac{2n}{k}$)};
    \node[segment, fill=gray!10, minimum width=2.4cm, anchor=west] (t5) at (t4.east) {$t_5$\\\footnotesize($\frac{2n}{k}$)};
    \node[draw=none, minimum width=1.2cm, anchor=west, font=\large\bfseries] (dots) at (t5.east) {$\dots$};
    \node[segment, fill=gray!10, minimum width=2.4cm, anchor=west] (tk) at (dots.east) {$t_{k/2}$\\\footnotesize($\frac{2n}{k}$)};

    % Imbalanced bucket Bi containing segment t3 entirely
    \node[bucketBox, minimum width=6.6cm, anchor=south] (bi) at ([yshift=14pt]t3.north) {
      Imbalanced Bucket $B_i$ with $|B_i| \ge \frac{4n}{k}$\\
      \footnotesize(Completely encloses $t_3 \implies t_3 \subseteq B_i$)
    };

    \draw[-{Stealth[scale=1.0]}, very thick, red!80!black] (bi.south) -- (t3.north);

    % Array endpoints
    \node[above=4pt, font=\scriptsize\bfseries] at (t1.north west) {$S_{\text{sort}}[1]$};
    \node[above=4pt, font=\scriptsize\bfseries] at (tk.north east) {$S_{\text{sort}}[n]$};
  \end{tikzpicture}%
  }
  \caption{Geometric containment construction: any bucket $B_i$ with size $\ge 4n/k$ spans across and fully encloses at least one segment $t^* \subset S_{\text{sort}}$ of width $2n/k$.}
  \label{fig:pigeonhole-chernoff-construction-en}
\end{figure}

\begin{lemma}[Geometric Containment Lemma]
\label{lem:containment-en}
If bucket $B_i$ contains $|B_i| \ge \frac{4n}{k}$ elements, it must \textbf{completely enclose at least one} canonical segment $t_j$:
\begin{equation}
  \exists j \in \left\{1, \dots, \frac{k}{2}\right\} \quad \text{such that} \quad t_j \subseteq B_i.
\end{equation}
\end{lemma}

\begin{proof}
Follows from the continuous pigeonhole principle. The separation boundaries between consecutive segments $t_j$ are spaced exactly $\frac{2n}{k}$ apart. An interval $B_i$ of length at least $\frac{4n}{k}$ has twice the width of a segment, thereby spanning across at least two consecutive boundaries and encompassing the full segment $t_j$ between them.
\end{proof}

\subsection{The Sample Count Contradiction}
Contrast this containment with the oversampling construction:
\begin{itemize}
  \item By definition, bucket $B_i$ is bounded by consecutive pivots $s_{i-1}$ and $s_i$.
  \item By pivot selection rule~\eqref{eq:pivot-selection-rule-en}, exactly $a$ samples from $A$ lie strictly between $s_{i-1}$ and $s_i$.
  \item Hence, $B_i$ contains at most $a$ samples from $A$.
  \item Because $t_j \subseteq B_i$, the number of samples from $A$ falling within $t_j$ cannot exceed the number in $B_i$.
\end{itemize}

This yields the fundamental deduction:
\begin{equation}
  \left( |B_i| \ge \frac{4n}{k} \right) \implies \left( \exists j \in \left\{1, \dots, \frac{k}{2}\right\} : |t_j \cap A| < a + 1 \right).
\end{equation}

Applying the Union Bound over all $\frac{k}{2}$ identical segments:
\begin{equation}
  \mathbb{P}(\text{failure}) \le \frac{k}{2} \cdot \mathbb{P}\left(t^* \text{ contains } < a + 1 \text{ samples from } A\right),
  \label{eq:union-bound-final-en}
\end{equation}
where $t^*$ represents an arbitrary fixed segment of width $\frac{2n}{k}$.

\subsection{Indicator Variables and Expectation}
Let $t^*$ be fixed. For each sample $A[r] \in A$ ($r = 1, \dots, |A|$), define the Bernoulli indicator:
\begin{equation}
  X_r \coloneqq \begin{cases}
    1 & \text{if sample } A[r] \in t^*, \\
    0 & \text{otherwise}.
  \end{cases}
\end{equation}

Because samples are chosen uniformly at random from $S$:
\begin{equation}
  p \coloneqq \mathbb{P}(X_r = 1) = \frac{|t^*|}{n} = \frac{\frac{2n}{k}}{n} = \frac{2}{k}.
\end{equation}

The total number of samples falling in $t^*$ is $X \coloneqq \sum_{r=1}^{|A|} X_r$. By linearity of expectation:
\begin{equation}
  \mathbb{E}[X] = |A| \cdot p = [(a + 1)k - 1] \cdot \frac{2}{k} = 2(a + 1) - \frac{2}{k}.
\end{equation}

Since $k \ge 2$, we have $\frac{2}{k} \le 1$. Given $a \ge 1$, we obtain:
\begin{equation}
  \mathbb{E}[X] \ge 2(a + 1) - 1 \ge \frac{3}{2}(a + 1).
  \label{eq:expected-value-bound-en}
\end{equation}

Inverting this relation yields:
\begin{equation}
  a + 1 \le \frac{2}{3} \mathbb{E}[X] = \left(1 - \frac{1}{3}\right) \mathbb{E}[X].
  \label{eq:delta-derivation-en}
\end{equation}

\subsection{Chernoff Lower Tail Bound}
The standard \textbf{Chernoff Lower Tail Bound}~\cite{chernoff1952,motwani1995} states that for independent Bernoulli trials and any deviation parameter $\delta \in (0, 1)$:
\begin{equation}
  \mathbb{P}(X < (1 - \delta)\mathbb{E}[X]) \le e^{-\frac{\delta^2}{2} \mathbb{E}[X]}.
  \label{eq:chernoff-lower-tail-en}
\end{equation}

Setting $\delta = \frac{1}{3}$:
\begin{equation}
  \mathbb{P}(X < a + 1) \le \mathbb{P}\left(X < \left(1 - \frac{1}{3}\right)\mathbb{E}[X]\right) \le e^{-\frac{(1/3)^2}{2} \mathbb{E}[X]} = e^{-\frac{1}{18} \mathbb{E}[X]}.
\end{equation}

Substituting $\mathbb{E}[X] \ge \frac{3}{2}(a + 1)$:
\begin{equation}
  -\frac{1}{18} \mathbb{E}[X] \le -\frac{1}{18} \cdot \frac{3}{2}(a + 1) = -\frac{a + 1}{12}.
\end{equation}

Inserting our choice $a = 12 \ln k$:
\begin{equation}
  -\frac{a + 1}{12} < -\frac{a}{12} = -\frac{12 \ln k}{12} = -\ln k.
\end{equation}

Exponentiating yields:
\begin{equation}
  \mathbb{P}(X < a + 1) \le e^{-\ln k} = \frac{1}{k}.
  \label{eq:single-interval-bound-en}
\end{equation}

\subsection{Proof Conclusion and Las Vegas Strategy}
Combining single-interval bound~\eqref{eq:single-interval-bound-en} with Union Bound~\eqref{eq:union-bound-final-en}:
\begin{equation}
  \mathbb{P}(\text{failure}) \le \frac{k}{2} \cdot \frac{1}{k} = \frac{1}{2}.
  \label{eq:total-failure-bound-en}
\end{equation}

\begin{theorem}[Success Guarantee of Over-Sampling]
With parameter $a = 12 \ln k$, the probability that \textbf{all} $k$ buckets are strictly balanced satisfies:
\begin{equation}
  \mathbb{P}(\text{success}) = 1 - \mathbb{P}(\text{failure}) \ge 1 - \frac{1}{2} = \frac{1}{2} \quad (50\%).
\end{equation}
\end{theorem}

\begin{property}[Zero Amortized I/O Las Vegas Restart]
If bucket validation detects an oversized bucket ($|B_i| \ge \frac{4n}{k}$), the algorithm discards the pivots and resamples independently. The number of trials $T \sim \text{Geom}(p)$ with $p \ge 1/2$ has expectation:
\begin{equation}
  \mathbb{E}[T] = \frac{1}{p} \le 2.
\end{equation}
Because sampling and sorting $A$ occur entirely in RAM with 0 disk I/O, retrying 1 or 2 times adds negligible overhead, preserving optimal $\mathcal{O}\left(\frac{n}{B} \log_{M/B}\frac{n}{B}\right)$ I/O performance.
\end{property}

\section{Uniform Random Sampling Algorithms}
\label{sec:random-sampling-algorithms-en}

The probabilistic guarantees of External Distribution Sort rely upon extracting truly uniform random samples without replacement.

\subsection{Problem Formulation}
Given a collection of $n$ elements $S[1 \dots n]$, extract a subset $D \subset S$ of fixed cardinality $|D| = m$ (here $m = |A| = (a+1)k - 1$) such that:
\begin{equation}
  \mathbb{P}(S[i] \in D) = \frac{m}{n}, \quad \forall i \in \{1, \dots, n\}.
\end{equation}

\subsection{Algorithm 1: Hash-Based Rejection Sampling}
The most direct approach iteratively generates pseudo-random indices in $[1, n]$, discarding duplicates via a hash set.

\begin{algorithm}[H]
\caption{Hash Set Random Sampling (\textit{HashSampling})}
\label{alg:hash-sampling-en}
\KwInput{Index range $[1 \dots n]$, sample size $m \le n$}
\KwOutput{Sorted set of $m$ distinct sampled indices $D$}

Initialize empty hash table $D \leftarrow \emptyset$\;
\While{$|D| < m$}{
  $p \leftarrow \text{rand}(1, n)$ \tcp*{Uniform integer in $[1, n]$}
  \If{$p \notin D$}{
    Insert $p$ into $D$\;
  }
}
Sort elements of $D$\;
\Return $D$\;
\end{algorithm}

\subsection{Collision Probability and Expected Trials}
When $|D| = j$, the probability that the next draw collides with an existing member is:
\begin{equation}
  \mathbb{P}(\text{collision at step } j) = \frac{j}{n}.
\end{equation}

Collision probability peaks on the final draw:
\begin{equation}
  \mathbb{P}(\text{collision}) \le \frac{m - 1}{n} < \frac{m}{n}.
\end{equation}

\begin{property}[Efficiency for Small Samples ($m \le n/2$)]
When $m \le \frac{n}{2}$ (as in oversampling where $m \ll n$), collision probability remains bounded by:
\begin{equation}
  \mathbb{P}(\text{collision}) < \frac{1}{2}.
\end{equation}
Thus, more than $50\%$ of generated numbers yield successful insertions.
\end{property}

The total expected number of random generations is:
\begin{equation}
  \mathbb{E}[\text{trials}] = \sum_{j=0}^{m-1} \frac{n}{n - j} = n \left( H_n - H_{n-m} \right) \approx n \ln\left( \frac{n}{n - m} \right).
  \label{eq:sampling-expected-attempts-en}
\end{equation}
Taylor expanding for $m \ll n$ gives $m + \frac{m^2}{2n} + \mathcal{O}(m^3/n^2)$. When $m = o(\sqrt{n})$, empty attempts are negligible, giving $\mathcal{O}(m)$ time.

\subsection{Structural Limitations and Access Models}
Despite simplicity, \autoref{alg:hash-sampling-en} has two architectural bottlenecks:

\begin{enumerate}
  \item \textbf{Variable-Length Keys and Strings:} Sampling index $p$ assumes $\mathcal{O}(1)$ random access. For variable-length strings or text records lacking fixed offsets, scalar byte positions cannot be indexed directly, requiring an auxiliary pointer index $P[1 \dots n]$ with $\Theta(n)$ space.
  
  \item \textbf{Absence of Streaming Property:} The algorithm cannot operate on data streams because it requires prior knowledge of $n$ and arbitrary random access across the dataset.
\end{enumerate}

\subsection{Streaming Paradigm: Reservoir Sampling}
When data arrives as an unbuffered stream with unknown length $n$, \textbf{Reservoir Sampling} (Algorithm R by Waterman and Vitter~\cite{vitter1985}) provides an optimal single-pass solution.

\begin{algorithm}[H]
\caption{Reservoir Sampling (\textit{Vitter's Algorithm R})}
\label{alg:reservoir-sampling-en}
\KwInput{Data stream $S = \langle x_1, x_2, \dots \rangle$, reservoir capacity $m$}
\KwOutput{Reservoir array $R[1 \dots m]$ containing $m$ uniform samples}

Initialize reservoir with first $m$ stream items: $R[1 \dots m] \leftarrow [x_1, \dots, x_m]$\;
$t \leftarrow m$\;
\While{stream yields next item $x_{t+1}$}{
  $t \leftarrow t + 1$\;
  $j \leftarrow \text{rand}(1, t)$ \tcp*{Uniform integer in $[1, t]$}
  \If{$j \le m$}{
    $R[j] \leftarrow x_t$ \tcp*{Replace element at index $j$}
  }
}
\Return $R$\;
\end{algorithm}

\begin{theorem}[Correctness of Reservoir Sampling]
After scanning any prefix of $n \ge m$ items from the stream, each item $x_i$ ($1 \le i \le n$) is present in reservoir $R$ with exact uniform probability:
\begin{equation}
  \mathbb{P}(x_i \in R) = \frac{m}{n}.
\end{equation}
\end{theorem}

\begin{proof}
By mathematical induction on stream length $n \ge m$.
\begin{itemize}
  \item \textbf{Base Case ($n = m$):} The first $m$ elements enter $R$ unconditionally. Probability is $m/m = 1$.
  \item \textbf{Inductive Step ($n > m$):} Assume that after $n - 1$ items, each item satisfies $\mathbb{P}(x_i \in R) = \frac{m}{n - 1}$.
  Upon the arrival of item $x_n$:
  \begin{enumerate}
    \item Item $x_n$ enters $R$ if $j \le m$, which occurs with probability $\frac{m}{n}$.
    \item A pre-existing element $x_i$ remains in $R$ if it is not replaced, which occurs with probability $1 - (\frac{m}{n} \cdot \frac{1}{m}) = \frac{n - 1}{n}$.
    \item By the induction hypothesis, $\mathbb{P}(x_i \in R \text{ at step } n) = \frac{m}{n - 1} \cdot \frac{n - 1}{n} = \frac{m}{n}$.
  \end{enumerate}
\end{itemize}
The theorem holds for all $n \ge m$.
\end{proof}

\section{Comparative Synthesis}
\label{sec:distribution-sort-synthesis-en}

\autoref{tab:sorting-paradigms-comparison-en} summarizes the trade-offs between external sorting paradigms, while \autoref{tab:sampling-strategies-comparison-en} contrasts uniform sampling strategies.

\begin{table}[H]
  \centering
  \caption{Architectural comparison between external sorting paradigms.}
  \label{tab:sorting-paradigms-comparison-en}
  \begin{tabularx}{\textwidth}{@{}>{\bfseries\arraybackslash}p{0.22\textwidth}p{0.36\textwidth}p{0.36\textwidth}@{}}
    \toprule
    Metric & Multiway Merge Sort (Chap. 3) & External Distribution Sort \\
    \midrule
    Philosophy & \textit{Bottom-Up} (progressive merging) & \textit{Top-Down} (recursive partitioning) \\
    \addlinespace
    I/O Complexity & $\mathcal{O}\left(\frac{N}{B} \log_{M/B} \frac{N}{B}\right)$ & $\mathcal{O}\left(\frac{N}{B} \log_{M/B} \frac{N}{B}\right)$ (high probability) \\
    \addlinespace
    RAM Buffers & $k$ read buffers, 1 write buffer & 1 read buffer, $k$ write buffers \\
    \addlinespace
    Final Phase & Complex $k$-way heap-based merge & Direct bucket concatenation (0 I/O) \\
    \addlinespace
    Vulnerability & High initial run generation traffic & Sensitivity to bucket skew \\
    \addlinespace
    Mitigation & Run size $2M$ via Loser Tree & Over-Sampling ($a = 12 \ln k$) with Chernoff \\
    \bottomrule
  \end{tabularx}
\end{table}

\begin{table}[H]
  \centering
  \caption{Analytical comparison of uniform random sampling algorithms.}
  \label{tab:sampling-strategies-comparison-en}
  \begin{tabularx}{\textwidth}{@{}>{\bfseries\arraybackslash}p{0.22\textwidth}p{0.22\textwidth}p{0.24\textwidth}Y@{}}
    \toprule
    Algorithm & Data Model & Time Cost & Characteristics and Limits \\
    \midrule
    Hash Table (\textit{Rejection}) & Random Access $\mathcal{O}(1)$ array & $\mathcal{O}(m)$ expected ($m \le n/2$) & Requires prior $n$; unusable on variable-length records without index. \\
    \addlinespace
    Reservoir Sampling (Vitter) & Data Stream (\textit{One-Pass}) & $\mathcal{O}(n)$ & No prior $n$ required; bounded $\mathcal{O}(m)$ RAM; strict sequential access. \\
    \addlinespace
    Partial Fisher-Yates & In-Place mutable array & $\mathcal{O}(m)$ & Mutates input array by swapping first $m$ entries; requires writable RAM. \\
    \bottomrule
  \end{tabularx}
\end{table}
